\documentclass[11pt,reqno]{amsart}
\usepackage[T1]{fontenc}
\usepackage[utf8]{inputenc}
\usepackage[english]{babel}
\usepackage{lmodern,microtype}
\usepackage{amsmath,amssymb,amsthm,mathtools,mathrsfs}
\usepackage{tikz-cd}
\usepackage{listings}
\usepackage[colorlinks=true,linkcolor=blue,citecolor=blue,urlcolor=blue]{hyperref}
\usepackage[nameinlink,noabbrev,capitalise]{cleveref}
\newtheoremstyle{mytheorem}{10pt}{10pt}{\itshape}{}{\bfseries}{.}{0.5em}{}
\newtheoremstyle{mydefinition}{10pt}{10pt}{\normalfont}{}{\bfseries}{.}{0.5em}{}
\theoremstyle{mytheorem}
\newtheorem{proposition}{Proposition}[section]
\theoremstyle{mydefinition}
\newtheorem{definition}[proposition]{Definition}
\newtheorem{remark}[proposition]{Remark}
\DeclareMathOperator{\Aut}{Aut}
\DeclareMathOperator{\hofib}{hofib}
\newcommand{\Type}{\mathsf{Type}}
\newcommand{\ag}[1]{\texttt{#1}}
\lstset{basicstyle=\ttfamily\small,frame=single,columns=fullflexible,keepspaces=true}

\title[Monodromy via Type Transport]{Monodromy via Type Transport:\\A Cubical Agda Formalization}
\author{Hongwei (Alex) Wang}
\address{New York University Abu Dhabi, Abu Dhabi, United Arab Emirates}
\email{hw4375@nyu.edu}
\date{2026-11-09}
\thanks{\textbf{Preliminary Draft --- Mathematical and Semantic Review Pending.} AI assistance was used in preparing the formalization and this source-informed reconstruction. It is not a mathematical author, and this attribution is not evidence of completed human review.}

\begin{document}
\begin{abstract}
We give a source-indexed Cubical Agda formalization of the constructions in a
frozen version of the nLab section ``Monodromy: In homotopy type theory.''
Starting from a family $P:B\to\Type$, the checked development identifies the
homotopy fiber of the dependent projection, constructs transport and its
equivalence structure, and records its unit and composition laws.  Restriction
to loops yields fiber monodromy.  For a chosen fiber point, this produces a
pointed equivalence to the transported basepoint and hence group equivalences
of based homotopy groups.  An explicit path between basepoints gives the
corresponding change-of-basepoint map.  No fundamental-group action or
higher-coherent action structure is claimed.
\end{abstract}
\maketitle

\section{Introduction}

This article accompanies the verified module \ag{NLabMonodromy.agda}.  Its
scope is deliberately fixed by the frozen nLab source recorded in the
experiment directory \cite{nlab}.  The source was available while this paper
was generated; hence this is a \emph{source-informed reconstruction}, not a
blind reverse-translation evaluation.  The purpose is to exhibit, declaration
by declaration, what the checked code expresses and what it does not express.

We use the standard perspective in which a dependent family is a fibration in
the homotopy-theoretic sense of type theory \cite[\S2.3]{hott}.  This does not
assert the classical Serre lifting property.  The relevant equality structure
is supplied by the type theory and made explicit below.

\section{Type families and homotopy fibers}

Let $B:\Type$ and $P:B\to\Type$.  Set
\begin{equation}\label{eq:total}
 E\coloneqq\sum_{b:B}P(b),\qquad \mathsf{pr}_1:E\to B.
\end{equation}
For $b:B$, the dependent projection occurs in the homotopy-pullback square
\[
\begin{tikzcd}[column sep=large,row sep=large]
 P(b) \arrow[r] \arrow[d] & E \arrow[d,"\mathsf{pr}_1"]\\
 {*} \arrow[r,"b"'] & B.
\end{tikzcd}
\]
The formal evidence is not a strict pullback assertion, but the checked
equivalence
\begin{equation}\label{eq:fiber}
 \hofib_{\mathsf{pr}_1}(b)\simeq P(b),
 \qquad \hofib_f(b)\coloneqq\sum_{e:E}(f(e)=b).
\end{equation}
In the module, the source-indexed declaration is
\ag{total-hofib}\ensuremath{\simeq}.  It is
obtained from the Cubical Library declaration \ag{fiberProjEquiv}, whose
orientation is reversed by the library operation \ag{invEquiv}.  The local
projection is exactly the following checked code excerpt.
\begin{lstlisting}
total-projection P = fst
\end{lstlisting}

\section{Transport and monodromy}

For $p:x=_B y$, transport in $P$ gives a function
\begin{equation}\label{eq:transport-function}
 p_*:P(x)\longrightarrow P(y).
\end{equation}
The local declaration \ag{transport-map} is intentionally kept distinct from
\ag{transport-equiv}, the latter being the equivalence $P(x)\simeq P(y)$.
The Cubical Library supplies the equivalence package through \ag{substEquiv};
its inverse data is mathematically transport along the inverse path.  The
actual local definition of the function is
\begin{lstlisting}
transport-map P p = subst P p
\end{lstlisting}

The two elementary compatibilities are checked as pointwise paths.  Thus for
$u:P(x)$,
\begin{equation}\label{eq:unit}
 (\mathsf{refl}_x)_*(u)=u,
\end{equation}
and, for $p:x=y$ and $q:y=z$,
\begin{equation}\label{eq:composition}
 (p\mathbin\cdot q)_*(u)=q_*\bigl(p_*(u)\bigr).
\end{equation}
Here Cubical Agda's $p\mathbin\cdot q$ means first $p$, then $q$.  Equation
\eqref{eq:composition} is represented by \ag{transport-comp}, using the
library path \ag{substComposite}; it is the evidence for the triangle
\[
\begin{tikzcd}[column sep=large,row sep=large]
 P(x) \arrow[r,"p_*"] \arrow[dr,"(p\cdot q)_*"'] & P(y) \arrow[d,"q_*"]\\
 & P(z).
\end{tikzcd}
\]
These are propositional, pointwise equalities.  They are not claims that the
two maps are definitionally equal.

Fixing $b_0:B$ and writing $F=P(b_0)$, transport along loops defines
\begin{equation}\label{eq:monodromy}
 \operatorname{mon}_{P,b_0}:\Omega(B,b_0)\longrightarrow\Aut(F),
 \qquad \gamma\longmapsto\gamma_*.
\end{equation}
The declaration \ag{monodromy-map} uses the library pointed loop space
\ag{Omega}; \ag{monodromy-refl} and \ag{monodromy-comp} provide the
corresponding pointwise unit and composition paths.  No higher coherence data
and no descent from the loop space to $\pi_1(B,b_0)$ is present.

\section{Pointed transport and homotopy groups}

Let $f_0:F$ and $gamma:b_0=b_0$.  Put $f_1\coloneqq\gamma_*(f_0)$.  The
important point is that $f_1=f_0$ is not a consequence of transport.  The
checked statement is instead the pointed equivalence
\begin{equation}\label{eq:pointed}
 (F,f_0)\simeq_*(F,f_1),
\end{equation}
represented by \ag{pointed-monodromy}.  Its basepoint witness is definitionally
the reflexivity path because the target is based at the transported point.

The library construction \ag{piIso} now yields a \emph{group} equivalence,
not merely an equivalence of carriers:
\begin{equation}\label{eq:pi}
 \pi_n(F,f_0)\cong\pi_n(F,f_1),\qquad n\geq1.
\end{equation}
The local declaration is \ag{pointed-monodromy-}\ensuremath{\pi}.  Cubical Library's
\ag{piGr n} presents conventional $\pi_{n+1}$, so the natural-number index
used in the code covers exactly the range $n\geq1$ in \eqref{eq:pi}.

Finally, if a path $r:f_1=f_0$ is supplied,
\ag{basepoint-change}\ensuremath{\bullet} and
\ag{basepoint-change-}\ensuremath{\pi} give the corresponding pointed
equivalence and group isomorphism.  The hypothesis $r$ remains explicit.  In particular, these
conditional maps do not supply an unconditional fundamental-group action.

\section{Formalization and verification}

The formal module has ten source statement identifiers, \ag{MON-HOTT-001}
through \ag{MON-HOTT-010}.  The total-space, transport, and monodromy
declarations are local source-indexed wrappers; the fiber, transport,
pointed-type, and homotopy-group mechanisms are attributed Cubical Library
results.  The pinned environment is Agda \ag{2.8.0-3d04bac} and Cubical Agda
Library v0.9, commit
\ag{87af036ebd1c87956cf08fea44e62c7cf76392ca}.

The module type-checks with the command recorded in the experiment report.
It has no holes, unsolved metas, local postulates, or termination overrides.
The JSONL alignment records give full Agda types, source locations, assumptions,
dependencies, and translation classifications for every frozen claim.

\section{Limitations}

Machine checking establishes that the formal declarations are accepted in the
pinned imported context.  It does not establish that this exposition and the
frozen source are semantically identical.  Mathematical correctness,
formalization correctness, and semantic alignment remain pending review by
Hongwei (Alex) Wang (REV-001).

The present code does not produce a fundamental-group action, show
representative-independence after truncation, choose basepoint-return paths,
or verify their compatibility with loop composition.  Those omissions are
intentional boundaries of the checked result, not suppressed proof gaps.

\begin{thebibliography}{9}
\bibitem{nlab}
nLab authors, \emph{Monodromy}, ``In homotopy type theory.''
\newblock \url{https://ncatlab.org/nlab/show/monodromy#InHomotopyTypeTheory}.

\bibitem{hott}
The Univalent Foundations Program, \emph{Homotopy Type Theory: Univalent
Foundations of Mathematics}, Institute for Advanced Study, 2013.
\newblock \url{https://homotopytypetheory.org/book/}.

\bibitem{cubical}
Cubical Agda Library contributors, \emph{Cubical Agda Library}, version 0.9,
commit \texttt{87af036ebd1c87956cf08fea44e62c7cf76392ca}.
\end{thebibliography}
\end{document}
